\chapter{Proof Construction}
\label{chp:proof}

In this chapter, we present the formal proof of the Isolation Lemma.
The proof follows the structure imposed by the formal development in Isabelle/HOL and proceeds by progressively reducing the failure of uniqueness to a simple probabilistic event.
Unlike pen-and-paper proofs, several steps that are implicit or informal in standard proofs must be mad explicit in formal settings, including existential witnesses, measuability, independence, type correctness, and totality.

In the following sections, we organize the proof as a sequence of structural reductions.
We begin by establishing the existence of minimizers and by reformulating its non-uniqueness in a constructive form.
Then we analyze the auxiliary quantity $\alpha(x)$, which plays a central role in isolating the dependence on individual weights.
Such analysis leads to a characterization of non-uniqueness as a collision event of the form $w(x) = \alpha(x)$.
Finally, we bound the probability of such events using the product probability space constructed in Chapter~\ref{chp:modeling}.

\section{Existence of Minimizers}
\label{sec:existence}

Before analyzing the uniqueness of minimizers, we must first ensure that the notion of a minimizer is well-defined.
In particular, for a given family of sets $\mathcal{F}$ and a weight function $w$, it is necessary to establish that at least one minimizer exists whenever $\mathcal{F}$ is non-empty.

Although the existence of a minimum is intuitively clear when weights take values in the natural numbers, such arguments cannot be taken for granted in the formal setting.
All existence claims must be derived explicitly from previously stated assumptions.
In a formal setting, this step is not merely a formality.
The following definitions and arguments rely on the this predicate as a non-vacuous object, and constructions such as the \texttt{LEAST} operator require the existence of witnesses to be established explicitly.
Isolating the existence of minimizers at this stage allows later proofs to proceed without repeatedly discharging existence conditions.

Formally, we show that for any non-empty family $\mathcal{F}$, there exists a set $A \in \mathcal{F}$ whose total weight is minimal among all elements of $\mathcal{F}$.
That is, $A$ satisfies the minimizer predicate introduced in Chapter~\ref{chp:modeling}.
This result is captured by the following lemma in Isabelle/HOL.

\medskip

\noindent
\textbf{Lemma (Existence of a Minimizer).}
\emph{If $\mathcal{F} \neq \emptyset$, then there exists a set
$A \in \mathcal{F}$ such that $A$ is a minimizer of $\mathcal{F}$ under
the weight function $w$.}

\medskip

This lemma is formalized as \texttt{exist\_minimizer} in the Isabelle development.
The proof constructs a concrete witness $A \in \mathcal{F}$ whose weight is minimal with respect to all other elements of $\mathcal{F}$, thereby ensuring that the minimizer predicate is non-vacuous.

The existence of such a minimizer allows us to reason about minimizers without introducing additional existence assumptions in subsequent arguments.
It also enables a constructive analysis of non-uniqueness, which will be the focus of the next section.

\section{Reformulating the Non-Uniqueness}
\label{sec:non-unique}

The Isolation Lemma concerns the probability that the minimizer of $\mathcal{F}$ under a random weight function $w$ is unique.
In order to reason about the failure of uniqueness in a formal setting, it is necessary to make this notion explicit and structurally usable.

In pen-and-paper proofs, non-uniqueness is often handeled informally by arguing that "there exist at least two different minimizers".
However, in a proof assistant such as Isabelle/HOL, statements involving negation of uniqueness cannot be used directly.
Instead, non-uniqueness must be reformulated as a positive existential statement that provides concrete witnesses.

From a logical perspective, this reformulation eliminates the negation of a unique existential statement to a explicit existential formulation.
Concretely, the minimizer of $\mathcal{F}$ under $w$ is not unique if and only if there exist two distinct sets $A$ and $B$ in $\mathcal{F}$ such that both satisfy the minimizer predicate.
That is, non-uniqueness can be charaterized by the existence of $A \neq B$ with
\[
A \in \mathcal{F}, \quad B \in \mathcal{F}, \quad
\text{and} \quad
\forall S \in \mathcal{F}, \;
w(A) \le w(S) \;\land\; w(B) \le w(S).
\]
Since both $A$ and $B$ satisfy the minimizer predicate, it follows immediately that they attain the same total weight.

This reformulation plays a crucial role in the formal proof.
By providing explicit witnesses $A$ and $B$, it allows subsequent arguments to fix such minimizers and analyze their structural relationship.
In particular, it enables reasoning about elements in the symmetric difference between $A$ and $B$, which forms the bases of the collision characterization developed later.

In the Isabelle/HOL development, this step is realized by unfolding the definition of the unique existential quantifier and reasoning with explicit minimizer predicates.
Although mathematically straightforward, making this reformulation explicit is essential for enabling later arguments that fix concrete minimizers and reason about their set-theoretic structure.


\section{The Alpha Construction}
\label{sec:alpha}

A central ingredient in the proof of the Isolation Lemma is the auxiliary quantity $\alpha(x)$ associated with each element $x \in U$.
Intuitively, $\alpha(x)$ isolates the dependence of the total weight of a set on the individual weight $w(x)$, allowing us to reason about the effect of modifying a single coordinate.

Formally, $\alpha(x)$ is defined as the difference between two minimal weight values determined by the family $\mathcal{F}$.
The first term represents the minimum total weight among all sets in $\mathcal{F}$ that do not contain $x$.
The second term represents the minimum residual weight required to form a set in $\mathcal{F}$ that contains $x$, excluding the contribution of $x$ itself.
This definition was introduced by Mulmuley et al.\ and plays a crucial role in isolating single-coorrdinate events.

\medskip

In our formalization, $\alpha(x)$ is defined using Isabelle's \texttt{LEAST} operator, which selects the least natural number satisfying a given predicate.
At this stage, $\alpha(x)$ is treated as an abstract quantity defined in terms of minimal weights, without assuming that the minima are realized by any particular sets.

\medskip

In the present chapter, we analyze its properties and show how the abstract \texttt{LEAST}-based definition can be simplifies in subsequent arguments, once the existence of minimizers has been derived.

\subsection{Simplifying the Definition of \texorpdfstring{$\alpha(x)$}{alpha(x)}}

The definition of $\alpha(x)$ involves two instances of the \texttt{LEAST} operator, each ranging over a family of sets in $\mathcal{F}$.
While this formalization is convenient at the modeling stage, it is too abstract to be used directly in the subsequent analysis, since the operator \texttt{LEAST} is not computationally usable in formal proof.

Once the existence of minimizers has been established, the two minimization terms appearing in the definition of $\alpha(x)$ can be simplified to concrete weighth values.
More precisely, it can be shown that the minimal values selected by the \texttt{LEAST} operator are realized by specific minimizers, depending on whether one element $x$ is contained in the minimizing set.

\medskip

First, suppose that $B$ is a minimizer of $\mathcal{F}$ under the weight function $w$, and that $x \notin B$.
In this case, $B$ satisfies the defining predicate of the first \texttt{LEAST} expression in the definition of $\alpha(x)$.
As a result, the minimum total weight among all sets in $\mathcal{F}$ that do not contain $x$ is attained by $B$, and the first \texttt{LEAST} expression can be shown equal to the concrete value $W(B)$.

\medskip

Second, suppose that $A$ is a minimizer of $\mathcal{F}$ and that $x \in A$.
Removing $x$ from $A$ yields a residual set whose total weight equals $w(A) - w(x)$.
By showing the $A$ satisfies the defining predicate of the second \texttt{LEAST} expression, the minimum total weight among all residual weights of sets in $\mathcal{F}$ that contain $x$ can be related to the concrete value $w(A \setminus \{x\})$.

\medskip

These observations are formalized by two auxiliary lemmas in the Isabelle development.
They connect the abstract \texttt{LEAST} expressions appearing in the definition of $\alpha(x)$ with concrete weights of minimizers.
Together, they allow $\alpha(x)$ to be rewritten in terms of ordinary arithmetic expressions whenever suitable minimizers are available.

\section{Independence of \texorpdfstring{$\alpha(x)$}{alpha(x)} from Individual Weights}
\label{sec:alpha_independence}

A crucial property of the auxiliary quantity $\alpha(x)$ is that it does not depend on the value of the individual weight $w(x)$.
Instead, $\alpha(x)$ is determined entirely by the weights assigned to those elements other than $x$.
This independence property forms the bridge between the structural analysis of minimizers and the probabilistic argument that follows.

To make this statement precise, consider two different weight functions $w$ and $w'$ that coincide on all elements of $u \setminus \{x\}$, but may differ in the value assigned to $x$.
Changing the value of $w(x)$ shifts the total weight of every set containing $x$ by the same offset, while leaving the relative ordering of weights determined by the remaining elements unchanged.
As a consequence, the minimal values appearing in the definition of $\alpha(x)$ can be shown to be invariant under such changes.

\medskip

Formally, we prove that if two weight functions agree on $U \setminus \{x\}$, then the corresponding values of $\alpha(x)$ are equal.
That is, $\alpha(x)$ depends only on the restrictions of the weight funcion to the complement of $\{x\}$.
This statement is captured by the following lemma in the Isabelle development.

\medskip

\noindent
\textbf{Lemma (Independence of $\alpha(x)$).}
\emph{Let $w$ and $w'$ be two weight functions such that $w(y) = w'(y)$ for all $y \in U \setminus \{x\}$.
Then $\alpha(x)$ computed with respect to $w$ is equal to $\alpha(x)$ computed with respect to $w'$.}

\medskip

This lemma is formalized as \texttt{alpha\_independent\_of\_coord} in Isabelle/HOL.
Its proof relies on the simplification of the \texttt{LEAST} expressions established in the previous section, together with basic congruence properties of addition and inequality on natural numbers.

The independence of $\alpha(x)$ from $w(x)$ is the key technical fact that enables the probabilistic analysis.
Once the weights of all other elements have been fixed, $\alpha(x)$ becomes a constant with respect to the remaining random choice of $w(x)$.
Consequently, the event $w(x) = \alpha(x)$ can be treated as a single-coordinate event in the product probability space, with a simple and uniform probability bound.

\section{Collision characterization via \texorpdfstring{$\alpha(x)$}{alpha(x)}}
\label{sec:collision}

We now establish the key structural link between the failure of uniqueness and the auxiliary quantity $\alpha(x)$.
Specifically, we show that if the minimizer of $\mathcal{F}$ under a weight function $w$ is not unique, then this non-uniqueness much be witnessed by a collision event of the form $w(x) = \alpha(x)$ for some element $x \in U$.

Suppose that the minimizer is not unique.
By the reformulation of non-uniqueness, this implies the existence of two distinct minimizer $A$ and $B$ in $\mathcal{F}$ under the same weight function $w$.
Since $A \neq B$, their symmetric difference in non-empty.
By symmetry, we may assume without loss of generality that there exists an element $x \in A \setminus B$.

\medskip

Because both $A$ and $B$ are minimizers, they have equal total weight, that is, $w(A) = w(B)$.
Applying the simplifications of the \texttt{LEAST} expressioons established in Section~\ref{sec:alpha}, we can express $\alpha(x)$ entirely in terms of the weight of $A$ and $B$.

Since $x \notin B$, the minimum total weight among all sets in $\mathcal{F}$ that do not contain $x$ is realized by $B$
Similarly, since $x \in A$, the minimum residual weight among all sets in $\mathcal{F}$ that contain $x$ in realized by the residual set $A \setminus \{x\}$.
Combining these observations yields
\[
\alpha(x)
= w(B) - w(A \setminus \{x\})
= w(A) - \bigl(w(A) - w(x)\bigr)
= w(x).
\]

Thus, the failure of uniqueness implies the existence of an element $x \in U$ such that the individual weight $w(x)$ coincides with $\alpha(x)$.

\medskip

This argument if formalized in Isabelle/HOL by the lemma \texttt{two\_minimizers\_alpha\_eq}, which shows that whenever two distinct minimizers exist, a corresponding collision event $w(x) = \alpha(x)$ must occur for some $x \in U$.

As a consequence, the event that the minimizer is not unique is contained in the union of the single-coordinate events $\{\, w(x) = \alpha(x) \mid x \in U \, \}$.
This reduction transforms the global failure of uniqueness into a collection of simple local events, paving the way for the probabilistic analysis inthe next section.

\section{Probabilistic Analysis and Final Bound}
\label{sec:probability}

We will now complete the proof of the Isolation Lemma by bounding the probability that the minimizer of $\mathcal{F}$ is not unique under a random weight assignment.
Throughout this section, weights are sampled according to the product probability space constructed in Chapter~\ref{chp:modeling}.
Unlike pen-and-paper arguments, all probabilistic statements here are interpreted with respect to an explicitly constructed probability mass function.

\medskip

By the collision characterization established in Section~\ref{sec:collision}, the event that the minimizer is not unique can be reduced to a union simple single-coordinate events.
More precisely, we have
\[
\{\, \text{the minimizer is not unique} \,\}
\subseteq
\bigcup_{x \in U} \{\, w(x) = \alpha(x) \,\}.
\]
This reduction transforms a global structural failure into a collection of local equality event, each involving only a single random variable.

\medskip

\paragraph{Single-coordinate probability bound.}
Fix an arbitrary element $x \in U$.
By the independence property proved in Section~\ref{sec:alpha_independence}, the auxiliary quantity $\alpha(x)$ depends only on the weights assigned to elements in $U \setminus \{x\}$.
Consequently, once the weights of all other elements have been fixed, $\alpha(x)$ becomes a constant with respect to the remaining random choice of $w(x)$.

Under the product probability space, the random variable $w(x)$ is distributed uniformly over the finite set $\{1 ,\dots,N\}$ and is independent of all other coordinates.
Consequently, the probability that $w(x)$ attains any particular value is exactly $1/N$.
In particular, we obtain the bound
\[
\Pr\bigl[\, w(x) = \alpha(x) \,\bigr] \le \frac{1}{N}.
\]

\medskip

\paragraph{Formalization considerations.}
Although the above argument is mathematically straightforward, its formalization requires several non-trivial steps.
In particular, the event $\{ w(x) = \alpha(x) \}$ must be shown to be a measurable event in the underlying probability space.
This relies on the measurability of the function $\alpha(x)$, which is itself defined using minimum operators over finite families of sets.

Furthermore, the independence arguments depend on the fact that the product probability space assigns weights to to different elements of $U$ independently.
In the Isabelle/HOL development, this property is derived explicitly from the construction of the product PMF using the \texttt{Pi\_pmf} operator, rather than being assumed implicitly.
These considerations account for a substantial portion of the formal proof effort, even though they are typically omitted in pen-and-paper proofs.

\medskip

\paragraph{Union bound and final estimate.}
Since the universe $U$ is finite, we may apply the union bound to the collection of collision events.
Summing the individual bounds yields
\[
\Pr\bigl[\, \text{the minimizer is not unique} \,\bigr]
\le \sum_{x \in U} \Pr\bigl[\, w(x) = \alpha(x) \,\bigr]
\le \frac{|U|}{N}.
\]
Equivalently, the probability that the minimizer is unique is bounded by
\[
\Pr\bigl[\, \text{the minimizer is unique} \,\bigr]
\ge 1- \frac{|U|}{N}
\]

\medskip

This inequality constitutes the statement of the Isolation Lemma.
In the Isabelle/HOL development, the above reasoning is formalized using standard results from the \texttt{HOL-Probability} library, including measurability lemmas for derived random variables and the union bound for finite collections of events.
The final result is captured by the theorem \texttt{isolation\_lemma\_main}, which completes the formal proof.


